\chapter{信息经济学与金融契约}
\label{chap:Information_Economics_and_Financial_Contracts}

\begin{openingcase}{为什么提高利率反而可能减少放贷？}

假设华信商业银行正在审查荣达机械和鑫海科技两家制造企业的贷款申请。从书面资料看，两家企业很相似：营业收入、杠杆、抵押品和近期财务表现大致相当，但它们的投资项目却截然不同。

荣达计划扩建一条成熟生产线，预期回报不高，但项目成功概率较大。鑫海则考虑一项投机性更强的技术投资：可能获得高得多的收益，但失败风险也明显更大。两家企业的管理者都远比银行更了解这些差异。

应对贷款风险的自然办法似乎是提高利率。贷款风险越高，银行就应要求更多补偿。然而，当借款人的差异无法被银行完全观察时，这一逻辑可能失效。

随着华信提高贷款利率，荣达这样的企业可能认为借款已不再划算。项目回报相对有限，难以承担高融资成本；鑫海这样的企业可能仍愿借款。如果项目成功，上行收益很大；如果失败，部分损失则由贷款人承担。

因此，利率影响的不只是贷款成功偿还时的承诺支付。它还可能影响\emph{谁愿意借款}，以及贷款发放后\emph{借款人愿意承担何种风险}。

当贷款利率足够高时，这些效应可能降低华信的预期回报。银行因而可能限制信贷供给，而不再继续加息，使一些愿意借款的企业无法获得贷款。

\end{openingcase}

\begin{tcolorbox}[
  enhanced,
  colback=FinanceBlue!6,
  colframe=FinanceBlue,
  boxrule=0.6pt,
  sharp corners,
  left=8pt,
  right=8pt,
  top=7pt,
  bottom=7pt,
  before skip=1em,
  after skip=1.6em,
  title=\textbf{核心问题},
  coltitle=white,
  colbacktitle=FinanceBlue,
  fonttitle=\bfseries,
  attach boxed title to top left={xshift=6pt,yshift=-2mm},
  boxed title style={
    boxrule=0pt,
    sharp corners,
    left=5pt,
    right=5pt,
    top=2pt,
    bottom=2pt
  }
]
\begin{enumerate}[
  leftmargin=*,
  label=\textcolor{FinanceBlue}{\textbf{\arabic*.}},
  itemsep=0.55em,
  topsep=0.2em
]
\item 签订契约之前，交易成本和信息不对称如何影响金融交换？

\item 提供融资之后，隐藏行动如何影响激励？

\item 契约、监督、抵押品、治理和执行如何缓解这些问题？为什么无法将其彻底消除？

\end{enumerate}
\end{tcolorbox}

金融契约往往在信息不对称条件下设计：借款人对自身前景和风险类型的了解，可能多于贷款人，而获得资金后的行为可能难以观察。签约前，私人信息影响借款人的选择；签约后，不可观察的行动可能改变风险承担激励。本章考察这些问题如何影响筛选、激励设计及金融风险分配。

\section{交易成本与信息摩擦}

金融交换具有成本。获取信息、评估交易对手、设计契约和执行契约，都需要资源。这些成本影响哪些交易能够达成、金融契约采用何种条款，以及金融关系如何组织。

\subsection{交易成本的组成}

考虑一家寻求外部融资的小型制造企业——明远精密零部件公司。提供资金之前，企业需要联系潜在贷款人或投资者，资金提供者则需要足够信息来评估机会。为此，可能要编制并审计财务报表、评估资产、起草法律文件和确定抵押品。各方还可能规定出现财务困境或违约时的权利、或有情形及救济措施。

上述活动都消耗资源。这些成本降低融资的净收益，并可能决定一笔交易能否产生正剩余。

假设明远考虑建设一条新生产线，今天需投资 $I$，下一期产生不确定支付 $X$。在无摩擦情形下，如果 $X$ 的预期现值超过初始投资，项目在财务上就值得实施。考虑交易成本 $K$ 后，相关剩余变为
\begin{equation}
S = V_{0}-I-K,
\label{eq:ch3-transaction-surplus}
\end{equation}
其中，$V_{0}$ 为项目预期未来支付的现值，$K$ 为安排和维持金融关系的交易成本现值。

\(K\) 的各组成部分发生在金融关系的不同阶段。寻找交易对手、评估相关信息时会产生搜寻和核实成本；契约成本来自对权利、义务及或有情形的规定。融资提供后，观察绩效与行为会产生监督成本；当契约义务未履行时，执行成本就变得重要。

许多活动涉及可以分摊到多笔交易上的固定成本，从而构成专业化的经济依据。家庭评估一笔企业贷款时，难以把信用分析和法律审查成本分摊到多笔交易上；相比之下，银行能够为众多借款人履行这些功能，将相关成本分摊到整个贷款业务中。审计师、信用评级机构和交易所，同样在核实、信息生产和标准化方面提供专业服务。

专业化不能消除交易成本，但可以降低每笔交易履行这些功能的平均成本。

\begin{figure}[H]
\centering
\begin{tikzpicture}[
    >=Latex,
    box/.style={
        draw,
        rounded corners,
        minimum width=2.7cm,
        minimum height=0.85cm,
        align=center
    },
    central/.style={
        draw,
        rounded corners,
        minimum width=3.2cm,
        minimum height=1.0cm,
        align=center,
        line width=0.9pt
    },
    flow/.style={
        -Latex,
        line width=0.8pt,
        FinanceBlue
    }
]

\node[central] (exchange) at (0,0) {金融交换};

\node[box] (search)   at (-4.3,1.6) {搜寻};
\node[box] (verify)   at (-4.3,0)   {核实};
\node[box] (contract) at (-4.3,-1.6) {缔约};

\node[box] (monitor) at (4.3,0.8) {监督};
\node[box] (enforce) at (4.3,-0.8) {执行};

\draw[flow] (search.east) -- (exchange.west);
\draw[flow] (verify.east) -- (exchange.west);
\draw[flow] (contract.east) -- (exchange.west);

\draw[flow] (exchange.east) -- (monitor.west);
\draw[flow] (exchange.east) -- (enforce.west);

\node[align=center] at (-4.3,2.45) {\small 签约之前及过程中};
\node[align=center] at (4.3,1.65) {\small 融资之后};

\end{tikzpicture}

\caption{金融关系中的交易成本}
\label{fig:ch3-transaction-costs}
\end{figure}

图~\ref{fig:ch3-transaction-costs} 是示意图，并不表示严格的时间顺序。签约后仍可能继续核实信息，监督可能产生新信息，而对执行的预期即使在违约发生前，也可能影响行为。

\subsection{信息不对称与有成本的信息获取}

当金融交易的一方掌握另一方不具备的相关信息时，就存在信息不对称。其中一些信息可能能够获得，但获取和核实需要成本。

这种不对称可能涉及特征。借款人对项目质量的了解可能多于贷款人，管理者对预期现金流的了解可能多于外部投资者。

它也可能涉及行动。贷款人可能无法完全观察借款人如何管理获融资项目；资本提供后，股东也可能观察不到重要的管理决策。

信息不对称不同于不确定性。即使所有人都能获得相同信息，未来结果仍可能不确定。区别在于信息在各方之间的分布，而非不确定性本身是否存在。

额外获取信息可以缩小部分信息差距，但需要消耗资源。令 \(V(q)\) 表示精度为 \(q\) 的信号所提供信息的经济价值，\(C(q)\) 表示生产该信号的成本。
假设决策者选择信息精度 \(q\)，以最大化净信息价值 \(V(q)-C(q)\)。在可微的内点最优解处，提高精度的边际价值等于边际成本：

\begin{equation}
V'(q)=C'(q).
\label{eq:ch3-information-acquisition}
\end{equation}

完全信息通常并非目标。关键是在给定成本下，获取多少信息值得。银行、投资者和保险公司因此会有选择地获取信息，在信息的预期价值与获取、处理所需资源之间权衡。

技术进步能够降低信息收集、处理和评估的成本，从而改变信息获取的经济性。数字记录、卫星数据、电子发票、机器学习和自动化文档分析，可以使过去获取成本高昂的信息更易获得，并更容易纳入金融决策。

然而，信息成本降低并不必然提高信息质量或决策价值。数据可能有偏、过时、被策略性操纵，或与关心的经济变量关联很弱。高效处理无法弥补信号经济相关性不足的问题。

\section{逆向选择：签约前的隐藏信息}

当私人掌握的特征影响参与或契约选择时，就会产生逆向选择（adverse selection）。不同类型对给定条款作出不同反应，而不知情一方在签约前无法完全观察每位参与者的类型。

\subsection{隐藏信息下的混同}

假设申请者中存在 $H$ 和 $L$ 两类借款人，$H$ 型项目质量高于 $L$ 型。贷款人知道申请者中 $H$ 型占比为 $\pi$，却无法在签约前观察任何单个申请人的类型。如果没有可以区分两类人的信息，就只能对申请者进行混同评估（pooling）。

若贷款人对两类人的估值分别为 $V_H$ 和 $V_L$，且 $V_H>V_L$，则类型未知的申请者的混同期望价值为

\begin{equation}
\bar V=\pi V_H+(1-\pi)V_L.
\label{eq:ch3-pooling-value}
\end{equation}

混同条件下，融资条款无法针对借款人的私人类型作出调整。如果参与决策取决于给定条款，高质量借款人可能更容易退出，使借款人池的平均质量下降。最初的混同估值便可能不再反映留下的借款人构成。这正是\emph{柠檬问题}（lemons problem）的核心机制。

这一机制出现在多种金融市场。如果贷款人无法区分借款人质量，低风险借款人可能面临部分反映高风险借款人风险的利率，从而减少外部融资。在保险市场，混同保费可能抑制低风险客户参与，却有利于高风险客户参与。在证券市场，发行人质量信息不足，可能导致高质量企业被低估。

\subsection{知情方的信号传递}

知情方可以通过可观察信号传递私人信息。只有当信号的成本或支付后果在不同类型之间存在系统性差异时，信号才有信息含量。

假设企业选择信号水平 $s$。类型 $j$ 的信号成本为 $c_j(s)$，其中 $j\in\{H,L\}$。单交叉条件可写为

\begin{equation}
c_H'(s)<c_L'(s),
\label{eq:ch3-single-crossing}
\end{equation}
即 $H$ 型传递信号的边际成本低于 $L$ 型。这一差异可以支持分离结果：高质量类型选择一种信号，而低质量类型认为模仿它并不值得。

金融信号可以包括保留更高股权比例、提供额外披露、接受更紧密地与绩效挂钩的报酬，或投入资产和个人资本。信号的信息含量，取决于不同类型在成本或承诺方面的差异。

能够以微不足道成本模仿的陈述，几乎不传递信息。如果低质量类型模仿的吸引力更小，付出成本的承诺就可能更可信。

信号传递（signaling）能够降低信息扭曲，但本身也有成本。用于区分类型的资源可能产生信息价值，却不会直接增加项目产出。

\subsection{不知情方的筛选}

筛选（screening）是指不知情方设计契约条款，诱使掌握私人信息的一方通过选择揭示信息。当不同类型偏好不同契约选项、其选择能够揭示私人信息时，这一机制就能发挥作用。

假设贷款人提供两种契约：契约 $A$ 利率较低，但抵押品要求较高；契约 $B$ 抵押品要求较低，但利率较高。不同类型借款人可能以不同方式评价这一权衡。因此，契约选择可以揭示贷款人无法直接观察的类型信息。

分离型契约菜单要求，每种类型都更偏好为其设计的契约。令 $C_H$ 和 $C_L$ 分别表示针对 $H$ 型和 $L$ 型的契约。若 $U_j(C_k)$ 表示类型 $j$ 选择契约 $C_k$ 的效用，其中 $j,k\in\{H,L\}$，则激励相容要求

\begin{equation}
\begin{aligned}
U_H(C_H) &\geq U_H(C_L),\\
U_L(C_L) &\geq U_L(C_H).
\end{aligned}
\label{eq:ch3-screening-ic}
\end{equation}

这些条件使自我选择与类型一致：$H$ 型至少弱偏好 $C_H$，$L$ 型至少弱偏好 $C_L$。当不同类型作出不同选择时，贷款人就能够从所选契约推断借款人类型。

基于契约的筛选在不同金融市场中形式各异。贷款人可以改变抵押品要求或贷款期限，保险公司则可以改变免赔额或保险限额。

如果不同类型偏好相似条款，其选择揭示的信息就较少。此时，可观察数据可以补充契约筛选，算法信用评分能够降低信息处理成本。但预测关系可能不稳定，也可能受到策略性操纵。

\subsection{市场后果：错误定价、信贷配给与市场瓦解}

在混同条件下，如果融资条款没有充分反映借款人质量差异，就可能出现错误定价。与完全信息相比，高质量借款人的条件可能更差，而低质量借款人的条件可能更好。

如果这些条款影响参与，一些高质量借款人可能退出。留下的借款人平均质量下降，进一步强化最初的定价问题。严重时，互利交易可能消失，市场部分或完全瓦解。

信贷配给（credit rationing）是贷款市场中尤其重要的后果。在标准竞争市场中，超额需求通常推动价格上涨。但在借款人质量隐藏的信贷市场，提高利率还可能改变愿意按该利率借款的借款人构成。

将贷款本金标准化为 1，假设非负项目支付可用于偿还贷款人，且不存在破产成本。令 \(\mathbb{E}_i\) 表示对愿意接受利率 \(i\) 的借款人求期望。银行每单位贷款的预期偿还额为

\begin{equation}
\rho(i)=\mathbb{E}_i[\min\{1+i,X\}]
\label{eq:ch3-bank-return}
\end{equation}
其中，$X$ 为每单位借款对应的项目支付。当 $X\geq 1+i$ 时，银行收到 $1+i$，否则收到 $X$。贷款利率变化同时影响契约偿还额和借款人类型的分布。


预期偿还额 $\rho(i)$ 未必随贷款利率单调增加。图~\ref{fig:ch3-credit-rationing} 展示了一种可能的非单调关系，其具体形式取决于借款人构成和项目支付如何随贷款利率变化。

\begin{figure}[H]
\centering

\begin{tikzpicture}[x=1cm,y=1cm]

% Axes
\draw[->, line width=0.8pt]
    (0,0) -- (7.2,0)
    node[right] {贷款利率 $i$};

\draw[->, line width=0.8pt]
    (0,0) -- (0,5.2)
    node[above] {预期偿还额 $\rho(i)$};

% Expected-return curve
\draw[line width=1.1pt, smooth]
    plot coordinates {
        (0.6,0.9)
        (1.4,1.8)
        (2.3,3.0)
        (3.2,4.0)
        (4.0,4.4)
        (4.8,4.2)
        (5.6,3.6)
        (6.5,2.6)
    };

% Maximum point
\fill (4.0,4.4) circle (2pt);

% Guide line
\draw[dashed, gray]
    (4.0,4.4) -- (4.0,0);

% Labels
\node[below] at (4.0,0) {$i^*$};
\node[above right] at (4.0,4.4)
    {\small 最大预期偿还额};

\end{tikzpicture}

\caption{贷款利率与预期偿还额之间的非单调关系示意}

\label{fig:ch3-credit-rationing}
\end{figure}


在图~\ref{fig:ch3-credit-rationing} 的示例中，令 $i^*$ 表示使预期偿还额最大化的贷款利率。若该利率下的贷款需求超过银行愿意提供的数量，银行可能选择配给信贷，而不再提高利率，因为进一步加息会降低预期偿还额。


\section{道德风险：签约后的隐藏行动}

与逆向选择不同，道德风险涉及签约后的隐藏行动。当这些行动影响另一方的支付，却无法被完全观察或写入可执行契约时，就可能出现激励问题。

\subsection{隐藏行动与代理问题}

隐藏行动是代理人在签约后选择的、影响委托人支付，却无法被完全观察或纳入契约的行动。当代理人的激励与委托人目标不完全一致时，这种信息限制就会产生代理问题（principal--agent problem）。

考虑提供资源的委托人和努力程度影响项目结果的代理人。项目成功时产生支付 $R$，失败时为零。代理人选择努力程度 $e$，它影响成功概率 $p(e)$，并产生私人成本 $c(e)$。

为简化分析，假设双方均为风险中性，因而货币支付线性进入其目标函数。若代理人在成功后获得 $w_S$，失败后获得 $w_F$，则预期效用为

\begin{equation}
U(e)=p(e)w_S+[1-p(e)]w_F-c(e).
\label{eq:ch3-agent-utility}
\end{equation}

委托人重视努力对项目回报的贡献，而代理人考虑契约报酬与努力的私人成本。如果努力不可观察，代理人只能获得其创造收益的一部分，便可能选择低于项目总剩余最大化水平的努力。

\subsection{努力、偷懒与激励问题}

有效率的努力水平 $e^*$ 使预期总剩余最大化：


\begin{equation}
e^*\in\arg\max_e\left\{p(e)R-c(e)\right\}.
\label{eq:ch3-efficient-effort}
\end{equation}

若 $e^*$ 是内点最优解，则满足

\begin{equation}
p'(e^*)R=c'(e^*).
\label{eq:ch3-efficient-effort-foc}
\end{equation}

代理人选择努力，以最大化契约提供的效用。将 $U(e)$ 对努力求导，内点最优解 $\hat e$ 满足

\begin{equation}
p'(\hat e)(w_S-w_F)=c'(\hat e).
\label{eq:ch3-agent-effort-foc}
\end{equation}

比较两个条件可识别激励问题的来源。项目从努力获得的边际收益为 $p'(e)R$，代理人的私人边际收益为 $p'(e)(w_S-w_F)$。若 $w_S-w_F<R$，代理人只能获得更高努力所创造增量价值的一部分。在关于 $p(e)$ 和 $c(e)$ 的标准假设下，私人选择的努力可能低于有效率水平。在此情形下，这一不足表现为偷懒（shirking）。

\begin{example}{高努力与低努力}{high-low-effort}

假设代理人在高努力 $e_H$ 与低努力 $e_L$ 之间选择。高努力提高成功概率，但私人成本也更高：

\begin{equation*}
p(e_H)>p(e_L),
\qquad
c(e_H)>c(e_L).
\end{equation*}

若代理人成功后获得 $w_S$、失败后获得 $w_F$，则在以下条件下更偏好高努力：

\begin{equation*}
p(e_H)w_S+[1-p(e_H)]w_F-c(e_H)
\geq
p(e_L)w_S+[1-p(e_L)]w_F-c(e_L).
\end{equation*}

整理得

\begin{equation*}
[p(e_H)-p(e_L)](w_S-w_F)
\geq
c(e_H)-c(e_L).
\label{eq:ch3-effort-spread}
\end{equation*}

左边是成功概率提高带来的额外预期报酬，右边是高努力的额外私人成本。成功与失败之间的报酬差越大，高努力相对于低努力就越有吸引力。

\end{example}

放松风险中性假设会引入另一项考虑。如果代理人厌恶风险，提高报酬对结果的敏感性，既强化激励，也使其在结果部分取决于不可控因素时面临更多收入风险。因此，契约设计可能需要在激励与风险分担之间权衡。

\subsection{风险转移与资产替代}

签约后的激励问题不仅来自努力选择，也来自项目选择。债务发行后，即使高风险项目降低企业总价值，股东仍可能更偏好它。这种激励称为风险转移（risk shifting）。当更安全的项目或资产被风险更高者替代时，通常称为资产替代（asset substitution）。

有限责任下，股东在满足承诺的债务偿还后获得剩余支付，其支付下限为零。如果企业支付为 $X$，承诺债务偿还额为 $D$，股东支付为

\begin{equation}
V_E(X)=\max\{X-D,0\}.
\label{eq:ch3-equity-payoff}
\end{equation}

这种不对称支付，可能使股东认为风险更高的项目更有吸引力，即使它降低企业总价值。

\begin{example}{债务发行后的风险转移}{risk-shifting-debt-issuance}

假设企业债务面值为 \(D=80\)。债务发行后，股东可在两个项目中选择。安全项目确定支付 110，股东获得 \(110-80=30\)。



风险项目以 \(0.5\) 的概率支付 150，以 \(0.5\) 的概率支付 50，其预期总支付为


\begin{equation*}
0.5(150)+0.5(50)=100,
\end{equation*}

低于安全项目的支付 110。在有限责任下，股东在高支付状态获得 70，低支付状态获得零。因此，其预期支付为



\begin{equation*}
0.5(70)+0.5(0)=35.
\label{eq:ch3-risk-shifting-example}
\end{equation*}

因此，尽管风险项目使企业总价值从 110 降至 100，股东仍偏好风险项目。债权人的预期偿还额则由安全项目下的 80 降为风险项目下的



\begin{equation*}
0.5(80)+0.5(50)=65
\end{equation*}

。股东价值增加的同时，债权人遭受更大损失，企业总价值下降。


\medskip

\begin{center}
\small
\textbf{两种项目选择下的支付}

\medskip

\begin{tabularx}{\linewidth}{
  @{}
  >{\raggedright\arraybackslash}X
  >{\centering\arraybackslash}p{2.8cm}
  >{\centering\arraybackslash}p{3.0cm}
  @{}
}
\toprule
 & \textbf{安全项目} & \textbf{风险项目} \\
\midrule
预期总支付
    & 110 & 100 \\
债权人预期偿还额
    & 80 & 65 \\
股东预期支付
    & 30 & 35 \\
企业价值排序 & 较高 & 较低 \\
股东支付排序 & 较低 & 较高 \\
\bottomrule
\end{tabularx}
\end{center}

\end{example}


融资后风险转移的可能性，为限制投资政策或经营风险变化的契约条款提供了依据。债权人关心的不仅是放贷时的条件，还有提供融资之后的决策。

\subsection{私人收益与代理成本}

当决策者获得的收益不按比例与外部投资者分享时，也可能出现签约后的激励冲突。管理者可能通过在职消费、建立企业帝国，或实施扩大自身控制却不增加企业价值的项目，获取私人收益。控股股东也可能以牺牲中小投资者利益为代价获益。

令决策 $a$ 产生企业价值 $V(a)$，并给决策者带来私人收益 $B(a)$。若决策者持有企业比例为 $\alpha$ 的股权，其私人支付为

\begin{equation}
V_{\text{私人}}(a)=\alpha V(a)+B(a).
\label{eq:ch3-private-benefits}
\end{equation}

当 $\alpha<1$ 时，决策者只承担企业价值下降的一部分，却保留全部私人收益。因此，使私人支付最大化的行动，可能不同于使企业总价值最大化的行动。

因此，私人收益可能使决策者私人支付最大化的行动，与企业价值最大化的行动之间产生偏离。这种偏离形成代理成本（agency costs），既包括用于限制激励冲突的资源，也包括冲突未能完全消除时留下的剩余损失。

\section{金融契约：应对信息与激励问题}

上述信息与激励问题，赋予契约设计以作用。金融契约不只是规定支付，还可以影响参与、揭示私人信息、协调激励，并在情形变化时分配权利。

披露、治理、审计、声誉和法律执行，往往支持这些功能。因此，金融契约是在更广泛的制度框架中运作的。

\subsection{抵押品与承诺}

契约设计能够将偿还义务与借款人未履约时仍可用的资产相联系。抵押品通过赋予债权人对指定资产的请求权，建立这种联系。借款人违约时，质押或抵押资产成为债权人的回收来源；违约之前，可能失去这些资产的后果也会影响借款人激励。

若欠款额为 $D$，抵押品可回收价值为 $C$，损失的简化度量为

\begin{equation}
L=\max\{D-C,0\}.
\label{eq:ch3-collateral-loss}
\end{equation}

抵押品还影响违约前的借款人激励。借款人面临失去抵押资产的可能，增加了提高违约风险行为的私人成本。因此，抵押品既是回收机制，也是承诺工具。

如果提供抵押资产的成本在借款人类型之间不同，抵押品还可能传递信息。自愿提供抵押可作为信号，抵押要求则可诱使借款人在契约之间自我选择。因此，其信息含量取决于不同类型提供抵押品的成本是否存在实质差异。

抵押品的有效性取决于抵押资产的价值和权利的可执行性。抵押品价值可能下降，扣押和变现可能代价高昂。抵押要求也可能限制有价值项目但可抵押资产不足的借款人融资。抵押品价值普遍下跌，会收紧借款约束并放大金融压力。

\subsection{限制性条款与控制权}

抵押品将偿还与特定资产相联系，限制性条款（covenants）则约束融资后的借款人行动。此类条款限制或要求特定行动，也可以规定影响债权人权利的条件。

消极性条款可能限制新增借款、资产出售或向股东分配；积极性条款可能要求披露财务信息、保全抵押品，或遵守特定财务比率要求。

限制性条款聚焦于有限的一组可观察或可核实的行动和财务条件。违反条款时，契约可能触发重新谈判，或向债权人转移更多控制权，例如要求补充抵押品、限制分配或加速偿还。

控制权往往取决于状态。满足契约条件时，管理者保留更多自主权；违反条件时，债权人获得更大影响力。

\begin{figure}[H]
\centering
\begin{tikzpicture}[
  >=Latex,
  box/.style={
    draw,
    rounded corners,
    minimum width=3.5cm,
    minimum height=0.85cm,
    align=center
  },
  flow/.style={
    -Latex,
    line width=0.85pt,
    FinanceBlue
  }
]

\node[box] (satisfied) at (-1,2.2)
  {满足契约条件};

\node[box] (manager) at (5.2,2.2)
  {管理者自主权较大};

\node[box] (breach) at (-1,0)
  {违反契约条件};

\node[box] (creditor) at (5.2,0)
  {债权人权利扩大};

\draw[flow] (satisfied) -- (manager);
\draw[flow] (breach) -- (creditor);

\draw[flow] (satisfied) --
  node[left, align=center]{\small 违反\\\small 限制性条款}
  (breach);

\end{tikzpicture}

\caption{通过限制性条款实现状态依存的控制权配置}
\label{fig:ch3-covenant-control}
\end{figure}

财务限制性条款需要在债权人及时介入与不必要的重新谈判之间权衡。门槛过严，可能因未反映财务状况实质恶化的暂时波动而触发违约；门槛过松，则可能在大量价值损失后才允许介入。

\subsection{支付设计与激励相容}

契约设计可以将支付与可观察结果联系起来，或提供鼓励自我选择的条款，以影响私人选择。激励相容（incentive compatibility）要求：筛选中每种类型都偏好针对自己的契约；行动隐藏时，所期望的行动对个人而言也是最优选择。

股权保留、绩效报酬和保险免赔额，能够通过提高个人对自身决策后果的承担来强化激励。然而，如果结果也取决于代理人无法控制的因素，更高的支付敏感性可能削弱风险分担。因此，契约设计需要在激励供给与风险分担之间权衡。

\subsection{监督与信息生产}

监督生产并核实融资后关于借款人绩效、合规和风险的信息。它可以揭示潜在违约，并支持契约权利的行使。行为可能被发现并受到执行的预期，也能抑制违反契约义务的行动。

当金融请求权高度分散时，监督可能不足。单个请求权人承担监督成本，却与其他请求权人共享收益，形成搭便车问题。金融中介可以代表多个资金提供者开展监督，从而缓解这一问题。

技术可以通过自动分析财务报告和交易数据来降低监督成本。AI 辅助系统能够识别潜在异常，但其结果仍需要核实和经济解释。

\subsection{治理与执行}

金融契约在更广泛的治理和执行安排中运行。治理关注管理者、董事会、股东、债权人和其他各方之间决策权与控制权的分配和行使。这些权利可以包括批准重大交易、更换管理者，或在特定条件满足时介入。

仅有权利分配并不保证其有效。执行决定契约和法律权利在实践中能否行使。法律规则和制度影响抵押品回收、加速偿还、争议解决，以及破产时债权优先顺序的落实。管理者或控股股东违反法定义务时，它们也提供救济。

执行不限于诉讼程序。声誉能够影响未来交易机会；承诺遭到违反时，贷款人和投资者可能撤回未来融资。交易所和监管机构也可以通过持续性要求和制裁来强化合规。

预期的执行有效性也影响契约设计。如果回收缓慢或不确定，贷款人可能要求更多抵押品、更短期限、更严条款或更高信用利差。更强的治理和更可靠的执行，可以在投资者较少直接介入的情况下，使外部融资成为可能。因此，同一份书面契约在不同制度环境中可能具有不同经济价值。

\begin{tcolorbox}[
  enhanced,
  breakable,
  colback=white,
  colframe=FinanceBlue,
  boxrule=0.7pt,
  sharp corners,
  borderline west={2.2pt}{0pt}{FinanceBlue},
  title=\textbf{金融视角：更好的信息并不能解决全部问题},
  coltitle=white,
  colbacktitle=FinanceBlue,
  fonttitle=\bfseries,
  before skip=1.2em,
  after skip=1.2em
]

更好的信息可以改善金融决策，却不能消除激励或执行问题。贷款人可能准确评估违约风险，却仍无法控制融资后的借款人行动。详细披露可以改善股东信息，却未必赋予外部投资者有效控制管理决策的权力。AI 模型可以识别信用质量恶化，但该信息的价值仍取决于可据此行使的权利和救济。因此，信息、激励和可执行权利是互补品，而非替代品。

\end{tcolorbox}

%========================================================
% Summary
%========================================================

\TailModule{summary}{本章小结}

金融交易需要有关交易对手、项目和未来行为的信息。生产和核实信息具有成本，而且相关信息往往在各方之间不对称分布。

逆向选择来自签约前的隐藏信息。质量无法直接观察时，混同定价可能抑制高质量参与者进入或留在市场，降低剩余参与者的平均质量，严重时促成市场瓦解。当不同类型模仿信号的成本存在差异时，信号传递使知情方能够揭示类型；筛选则利用契约选择诱导自我选择。在信贷市场，提高利率可能恶化借款人构成，因此银行预期偿还额未必随利率单调增加。即使有借款人愿意接受更高利率，也可能出现信贷配给。

道德风险来自签约后的隐藏行动。偷懒、风险转移和追求私人收益，会使私人最优行动偏离贷款人、投资者或保险公司所偏好的行动。金融契约通过改变支付、承诺和控制权予以应对。抵押品、限制性条款、股权保留、免赔额、监督及状态依存控制可以改善激励，而治理和执行决定权利在实践中如何行使。

这些机制都涉及权衡。强化激励可能削弱风险分担，更严的契约限制可能压缩有价值的灵活性，抵押要求可能限制可抵押资产较少的借款人，而监督本身也有成本。因此，契约设计旨在以合理经济成本减少信息和激励问题，而非彻底消除它们。

%========================================================
% Mind Map
%========================================================

\TailModule{mindmap}{思维导图}
\begin{figure}[H]
\centering
\begin{tikzpicture}[
  >=Latex,
  main/.style={
    draw,
    rounded corners,
    minimum width=3.4cm,
    minimum height=0.9cm,
    align=center,
    fill=gray!10
  },
  sub/.style={
    draw,
    rounded corners,
    minimum width=3.2cm,
    minimum height=0.8cm,
    align=center
  },
  flow/.style={-Latex, line width=0.75pt}
]

\node[main] (friction) at (0,4.2)
  {信息摩擦};

\node[sub] (hiddeninfo) at (-4.0,2.3)
  {签约前的\\隐藏信息};

\node[sub] (hiddenaction) at (4.0,2.3)
  {签约后的\\隐藏行动};

\node[sub] (adverse) at (-4.0,0.5)
  {逆向选择};

\node[sub] (moral) at (4.0,0.5)
  {道德风险};

\node[sub] (responses1) at (-4.0,-1.5)
  {信号传递与筛选};

\node[sub] (responses2) at (4.0,-1.5)
  {支付设计与监督};

\node[main] (responses) at (0,-3.4)
  {契约与\\制度应对};

\draw[flow] (friction) -- (hiddeninfo);
\draw[flow] (friction) -- (hiddenaction);
\draw[flow] (hiddeninfo) -- (adverse);
\draw[flow] (hiddenaction) -- (moral);
\draw[flow] (adverse) -- (responses1);
\draw[flow] (moral) -- (responses2);
\draw[flow] (responses1) -- (responses);
\draw[flow] (responses2) -- (responses);

\end{tikzpicture}

\caption*{签约前后的信息问题}
\end{figure}

%========================================================
% Key Terms
%========================================================

\TailModule{keyterms}{关键术语}

交易成本（Transaction costs）；信息不对称（information asymmetry）；有成本的信息获取（costly information acquisition）；隐藏信息（hidden information）；逆向选择（adverse selection）；混同（pooling）；信号传递（signaling）；筛选（screening）；激励相容（incentive compatibility）；信贷配给（credit rationing）；隐藏行动（hidden action）；道德风险（moral hazard）；委托代理问题（principal--agent problem）；风险转移（risk shifting）；资产替代（asset substitution）；私人收益（private benefits）；代理成本（agency costs）；抵押品（collateral）；限制性条款（covenant）；控制权（control rights）；监督（monitoring）；治理（governance）；执行。


%========================================================
% Financial AI Lab
%========================================================（enforcement


%========================================================
% Financial AI Lab
%========================================================）。

\TailModule{aiprompt}{金融 AI 实验室}

\textbf{实验：设计并审查 AI 信贷筛选框架}

请 AI 系统为一家小企业贷款机构提出信贷筛选框架，要求识别可观察变量，并解释每个变量可能如何关联偿还风险。区分放贷前可获得的信息，与融资后才能获得的信息，并解释后者为什么属于监督而非初始筛选。

随后审查框架：哪些变量包含借款人内在质量的信息，哪些只是与之相关，哪些会在申请人理解筛选规则后被操纵？指出一种潜在的选择偏差来源、一项公平性问题，以及一个使用历史借款人训练的模型可能在信贷条件变化后表现不佳的原因。

最后，假设模型能够高度准确地预测违约风险。解释哪些问题仍需通过契约条款、抵押品、监督、执行或人工审查来解决。为什么准确预测本身并不能消除缔约问题？

%========================================================
% Discussion Questions
%========================================================

\TailModule{questions}{讨论题}

\begin{enumerate}

\item 为什么不确定性不等同于信息不对称？分别举出一个不确定性很高但信息对称的金融例子，以及一个信息不对称的例子。

\item 为什么提高贷款利率，可能在承诺偿还额上升的同时，降低银行每单位贷款的预期偿还额？


\item 比较信号传递与筛选。两种情况下分别由谁主动采取措施？什么机制使不同类型能够被区分？

\item 为什么融资后对借款人行为的信息改善，未必能消除道德风险？解释监督、契约激励、控制权和执行的作用。



\end{enumerate}

%========================================================
% Applied Problems
%========================================================

\TailModule{appprobs}{应用题}

\begin{enumerate}

\item 一位投资者考虑购买对企业未来项目现金流的请求权。他认为其中 70\% 为高质量，每项价值 120；30\% 为低质量，每项价值 60。他无法观察单项请求权的质量。
    \begin{enumerate}

    \item 计算混同期望价值。
    
    \item 假设高质量企业只愿以至少 110 的价格出售，低质量企业只要求至少 55。如果投资者最初按混同期望价值报价，哪些类型会接受？出售请求权的构成如何变化？

    \item 从经济角度解释逆向选择机制。

    \end{enumerate}

\item 一名代理人在高、低努力之间选择。高努力使成功概率为 $0.8$，私人成本为 12；低努力使成功概率为 $0.5$，私人成本为 3。代理人成功后获得 $w_S$，失败后获得零。

  \begin{enumerate}

  \item 推导选择高努力的激励相容条件，并确定 $w_S$ 的门槛值。

  \item 解释为什么提高成功时的支付会影响代理人的努力选择。

  \item 如果代理人厌恶风险，会出现什么额外权衡？

  \end{enumerate}

\item 某企业债务面值为 80。股东可以选择项目 A，确定支付 112；或项目 B，以 $0.5$ 的概率支付 160，以 $0.5$ 的概率支付 48。

  \begin{enumerate}

  \item 计算每个项目的预期总支付。

  \item 计算每个项目下股东的预期支付。

  \item 股东偏好的项目是否使企业总价值最大化？解释冲突的来源。

  \end{enumerate}

\end{enumerate}

%========================================================
% Case Analysis
%========================================================

\TailModule{caseanalysis}{案例分析}

\textbf{瑞幸咖啡：当披露信息不能照单全收}

2020 年 12 月，美国证券交易委员会（SEC）指控瑞幸咖啡公司通过严重错报收入、费用和净营业亏损欺诈投资者。根据 SEC 诉状，至少从 2019 年 4 月至 2020 年 1 月，瑞幸故意虚构了超过 3 亿美元的零售销售额。公司在不承认也不否认指控的情况下，同意一项有待法院批准的和解，其中包括支付 1.8 亿美元民事罚款。

以 SEC 于 2020 年 12 月 16 日发布的执法公告及相关公司披露为一手来源，运用本章框架分析这一事件。外部投资者依赖哪些信息？哪些信息生产和核实机制原本用于减少内部人与外部投资者的信息不对称？此事件反映了外部审计作为核实机制的哪些价值和局限？披露信息受到质疑后，哪些监督、治理、内部控制和执行机制变得重要？最后，区分案例中的\emph{信息失灵}与\emph{激励或治理失灵}。

\emph{一手来源：美国证券交易委员会，“瑞幸咖啡同意支付 1.8 亿美元罚款以和解会计欺诈指控”（``Luckin Coffee Agrees to Pay \$180 Million Penalty to Settle Accounting Fraud Charges''），新闻稿 2020-319，2020 年 12 月 16 日。}
%========================================================
% Further Reading
%========================================================

\TailModule{furtherread}{延伸阅读}

\begin{itemize}

\item Akerlof, George A.（1970），“柠檬市场：质量不确定性与市场机制”（``The Market for `Lemons': Quality Uncertainty and the Market Mechanism''），\emph{The Quarterly Journal of Economics}，84(3)，488--500。

\item Spence, Michael（1973），“劳动力市场信号传递”（``Job Market Signaling''），\emph{The Quarterly Journal of Economics}，87(3)，355--374。

\item Jensen, Michael C. 与 William H. Meckling（1976），“企业理论：管理者行为、代理成本与所有权结构”（``Theory of the Firm: Managerial Behavior, Agency Costs and Ownership Structure''），\emph{Journal of Financial Economics}，3(4)，305--360。

\item Stiglitz, Joseph E. 与 Andrew Weiss（1981），“信息不完全市场中的信贷配给”（``Credit Rationing in Markets with Imperfect Information''），\emph{American Economic Review}，71(3)，393--410。

\item Diamond, Douglas W.（1984），“金融中介与委托监督”（``Financial Intermediation and Delegated Monitoring''），\emph{The Review of Economic Studies}，51(3)，393--414。

\end{itemize}
\ChapterTail{summary,mindmap,keyterms,aiprompt,questions,appprobs,caseanalysis,furtherread}
